\documentclass{article}
\usepackage[english]{babel}

\usepackage[a4paper,top=2cm,bottom=2cm,left=3cm,right=3cm,marginparwidth=1.75cm]{geometry}
\usepackage{amsmath}
\usepackage{unicode-math}
\usepackage{graphicx}
\usepackage[colorlinks=true, allcolors=blue]{hyperref}
\usepackage{listings}
\usepackage{minted}
\usepackage{booktabs}
\usepackage{physics}
\usepackage{amssymb}%braket schreibweise
\setlength{\parindent}{0pt}

\title{Physics of Complex Systems: Exercise }
\author{Leopold Schaller and Daniel Lender}

\begin{document}
\maketitle

\section*{Task 8.1: Fireflies}

    \subsection*{(a)}
    We need an output range within the interval $[0,1]$.
    The activation function in the last layer must give an output accordingly. This requirement is fulfilled by the logistic function 
    \begin{align}
        \sigma(x) = \frac{1}{1 + e^{-\beta x}}
    \end{align}
    which gives an output in exactly this interval.

    \subsection*{(b)}
    Given that the probability of a firefly emitting light is $\langle t_i\rangle=p_i$, the complementary event is $\langle1- t_i\rangle=1-p_i$. As the expectation operator is linear, we can substitute these expected values into the cross-entropy formula:
    \begin{align}
    H= \langle H_S \rangle=−\frac{1}{N}\sum^N_{i=1}\left(\langle t_i\rangle \text{log}_2(q_i)+(\langle 1− t_i\rangle)\text{log}_2(1−q_i)\right)
    \end{align}
    
    Substituting $\langle t_i\rangle=p_i$ yields the final expectation value:
    \begin{align}
        H =−\frac{1}{N}\sum^N_{i=1}(p_i\text{log}_2(q_i)+(1−p_i)\text{log}_2(1−q_i))
    \end{align}

    \subsection*{(c)}
    By the Central Limit Theorem, the variance of the sample mean decreases inversely with the dataset size:\begin{align}
        \text{Var}(H_S) \propto \frac{1}{N}
    \end{align}
    For a sufficiently large dataset, the statistical fluctuations around the expected value become negligibly small, making $H_S \approx \langle H_S \rangle = H$ a mathematically robust assumption.

    \subsection*{(d)}
    Rewrite the objective function using the natural logarithm ($\log_2(x) = \frac{\ln(x)}{\ln(2)}$):
    \begin{align}
        H = -\frac{1}{N \ln 2} \sum_{i=1}^N \Big( p_i \ln(q_i) + (1 - p_i) \ln(1 - q_i) \Big)
    \end{align}

    Verify $q_i = p_i$ is a critical point:
    \begin{align}
        \frac{\partial H}{\partial q_i} = -\frac{1}{N \ln 2} \left( \frac{p_i}{q_i} - \frac{1 - p_i}{1 - q_i} \right) = 0
    \end{align}
    This is only satisfied if 
    \begin{align}
        \frac{p_i}{q_i}&=\frac{1-p_i}{1-q_i}\\[0.2cm]
        p_i(1 - q_i) &= q_i(1 - p_i)\\[0.2cm]
        p_i - p_i q_i &= q_i - p_i q_i\\[0.2cm]
       \Rightarrow q_i &= p_i
    \end{align}

    Verify the critical point at $q_i = p_i$ is the global minimum:
    \begin{align}
        \frac{\partial^2 H}{\partial q_i^2} 
        &= -\frac{1}{N \ln 2} \left( -\frac{p_i}{q_i^2} - \frac{1 - p_i}{(1 - q_i)^2} \right)\\
        &= \frac{1}{N \ln 2} \left( \frac{p_i}{q_i^2} + \frac{1 - p_i}{(1 - q_i)^2} \right)
    \end{align}
    Since $p_i$ and $q_i$ lie in the interval $[0,1]$, the second derivative is strictly positive. This means, H is strictly convex, i.e., proves the critical point at $q_i = p_i$ is the global minimum.
    

    \subsection*{(e)}
    The general Kullback-Leibler divergence is
    \begin{align}
        D_{KL}(P||Q)= \sum_{x\in\mathcal{X}} P(x)\ln\left(\frac{P(x)}{Q(x)}\right)
        \intertext{for a binary distribution}
        D_{KL, binary}(p \parallel q) = p \log_2 \frac{p}{q} + (1-p) \log_2 \frac{1-p}{1-q}
    \end{align}
% https://doi.org/10.1002%2F047174882X  chapter 2.4
    Calculating $H-H_{min}$: 
    \begin{align}
        H -H_{min} &= -\frac{1}{N \ln 2} \left[\sum_{i=1}^N \Big( p_i \ln(q_i) + (1 - p_i) \ln(1 - q_i) \Big) - \sum_{i=1}^N \Big( p_i \ln(p_i) + (1 - p_i) \ln(1 - p_i) \Big) \right]\\
        &= -\frac{1}{N \ln 2} \sum_{i=1}^N \Big( p_i \ln(q_i) + (1 - p_i) \ln(1 - q_i)  -   p_i \ln(p_i) - (1 - p_i) \ln(1 - p_i) \Big) \\
        &=-\frac{1}{N \ln 2} \sum_{i=1}^N \Big( p_i \big[\ln(q_i)-\ln(p_i)\big] + (1 - p_i) \big[\ln(1 - q_i)  -    \ln(1 - p_i)\big] \Big)\\
        &=-\frac{1}{N \ln 2}\sum_{i=1}^N \Bigg( p_i \ln\left(\frac{q_i}{p_i}\right) + (1 - p_i)\ln\left(\frac{1 - q_i}{1-p_i}\right) \Bigg)\\
        &=\frac{1}{N}\sum_{i=1}^N \Bigg( p_i \log_2\left(\frac{p_i}{q_i}\right) + (1 - p_i)\log_2\left(\frac{1 - p_i}{1-q_i}\right) \Bigg)\\
        &=\frac{1}{N}\sum_{i=1}^N  D_{KL, binary}(p \parallel q) 
    \end{align}
    Therefore, the difference $H - H_{min}$ is the average Kullback-Leibler divergence across all $N$ fireflies.



\section*{Task 8.2: Cross Entropy}

    \subsection*{(a)}
    \begin{align}
        \mathcal{E}= -\sum_{\mu=1}^M p_\mu(\Vec{x}) \ln{q_\mu(\Vec{x})}
    \end{align}

    \subsection*{(b)}
    Two constraints have to be taken into account:
    \begin{align}
    \sum_{\mu=1}^M q_\mu(\Vec{x})=1\hspace{2cm} \text{and} \hspace{2cm}
     q_\mu \in [0,1]
    \end{align}
    The first constraint makes sure, that the predicted probabilities for all outcomes sums to one, i.e., the distribution is normalized.
    The second, that all $q_\mu$ are actual probabilities (non negative and not greater than one).
    
    \subsection*{(c)}
    The general Lagrange function:
    \begin{align}
        \mathcal{L}=f(x)+\langle\lambda, g(x)\rangle
        \intertext{where $\lambda$ is the Lagrange multiplier. For our case we use the dot product for the inner product and consider a optimization problem with one constraint:
        \begin{itemize}
            \item minimize $f(x)$
            \item subject to $g(x)=0$
        \end{itemize}
        Our Lagrangian with the multiplier $\lambda$ is then}
        \mathcal{L}(q, \lambda) = -\sum_{\mu=1}^{M} p_\mu \log q_\mu + \lambda \left( \sum_{\mu=1}^{M} q_\mu - 1 \right)
    \end{align}
    To find the extremum, we take the partial derivative of $\mathcal{L}$ with respect to a specific predicted probability $q_\mu$ and set it to zero:
    \begin{align}
    \frac{\partial \mathcal{L}}{\partial q_\mu} &= -\frac{p_\mu}{q_\mu} + \lambda = 0\\[0.2cm]
    \Rightarrow q_\mu&=\frac{p_\mu}{\lambda}
    \intertext{now sum both sides over all possible outcomes}
    \sum_{\mu=1}^{M} q_\mu = \sum_{\mu=1}^{M} \frac{p_\mu}{\lambda}
    \intertext{now we apply the constraint $\sum_{\mu=1}^{M} q_\mu = 1$:}
    1= \sum_{\mu=1}^{M} \frac{p_\mu}{\lambda}
    \intertext{Since $p$ is a true probability distribution, its probabilities also sum to 1 ($\sum_{\mu=1}^{M} p_\mu = 1$):}
    1 = \frac{1}{\lambda} (1) \Rightarrow \lambda = 1
    \end{align}
    Substituting $\lambda = 1$ back into the expression for $q_\mu$:
    \begin{align}
        q_\mu = p_\mu \quad \forall \mu \in \{1 \dots M\}
    \end{align}


\end{document}

# lualatex -shell-escape Uebung6/exercise6/exercise6_Schaller_Lender.tex
